\documentclass[11pt,a4paper]{article}

\usepackage[utf8]{inputenc}
\usepackage[T1]{fontenc}
\usepackage{lmodern}
\usepackage[french]{babel}

\usepackage[margin=2.5cm]{geometry}
\usepackage{amsmath,amssymb}
\usepackage{xcolor}
\usepackage{booktabs}
\usepackage{hyperref}
\usepackage{tcolorbox}
\tcbuselibrary{theorems,breakable,skins}
\usepackage{enumitem}
\usepackage{fancyhdr}
\usepackage{array}
\usepackage{pifont}
\usepackage{listings}

\newcommand{\cmark}{\ding{51}}

\hypersetup{
    colorlinks=true,
    linkcolor=blue!60!black,
    urlcolor=blue!60!black,
    citecolor=blue!60!black
}

\newtcolorbox{objectifbox}[1][]{colback=blue!5,colframe=blue!50,
    fonttitle=\bfseries,title=#1,breakable}
\newtcolorbox{infobox}[1][]{colback=green!5,colframe=green!50!black,
    fonttitle=\bfseries,title=#1,breakable}
\newtcolorbox{attentionbox}[1][]{colback=orange!5,colframe=orange!70,
    fonttitle=\bfseries,title=#1,breakable}
\newtcolorbox{retenirbox}[1][]{colback=yellow!10,colframe=yellow!50!black,
    fonttitle=\bfseries,title=#1,breakable}
\newtcolorbox{prereqbox}[1][]{colback=cyan!5,colframe=cyan!50!black,
    fonttitle=\bfseries,title=#1,breakable}
\newtcolorbox{exercicebox}[1][]{colback=purple!5,colframe=purple!50,
    fonttitle=\bfseries,title=#1,breakable}

\lstdefinestyle{java}{
  language=Java,
  basicstyle=\ttfamily\small,
  backgroundcolor=\color{gray!8},
  frame=single, framerule=0.5pt, rulecolor=\color{gray!40},
  keywordstyle=\color{blue!70!black}\bfseries,
  commentstyle=\color{green!50!black}\itshape,
  stringstyle=\color{red!60!black},
  numbers=left, numberstyle=\tiny\color{gray},
  breaklines=true, showstringspaces=false, tabsize=4,
  xleftmargin=4pt, xrightmargin=4pt, aboveskip=8pt, belowskip=8pt,
  literate={é}{{\'e}}1 {è}{{\`e}}1 {ê}{{\^e}}1 {à}{{\`a}}1 {ù}{{\`u}}1
           {ô}{{\^o}}1 {î}{{\^i}}1 {ç}{{\c{c}}}1 {É}{{\'E}}1
}

\lstdefinestyle{bash}{
  language=bash,
  basicstyle=\ttfamily\small,
  backgroundcolor=\color{black!5},
  frame=single, framerule=0.5pt, rulecolor=\color{gray!40},
  commentstyle=\color{green!50!black}\itshape,
  numbers=none, breaklines=true, showstringspaces=false,
  xleftmargin=4pt, xrightmargin=4pt, aboveskip=8pt, belowskip=8pt,
  % Le style « bash » était le seul des trois sans table literate : un « é » dans un bloc bash
  % — « # Se déconnecter » — faisait échouer la compilation. Même table que java et python.
  literate={é}{{\'e}}1 {è}{{\`e}}1 {ê}{{\^e}}1 {à}{{\`a}}1 {ù}{{\`u}}1
           {ô}{{\^o}}1 {î}{{\^i}}1 {ç}{{\c{c}}}1 {É}{{\'E}}1
}

\lstdefinestyle{python}{
  language=Python,
  basicstyle=\ttfamily\small,
  backgroundcolor=\color{yellow!5},
  frame=single, framerule=0.5pt, rulecolor=\color{gray!40},
  keywordstyle=\color{blue!70!black}\bfseries,
  commentstyle=\color{green!50!black}\itshape,
  stringstyle=\color{red!60!black},
  numbers=left, numberstyle=\tiny\color{gray},
  breaklines=true, showstringspaces=false, tabsize=4,
  xleftmargin=4pt, xrightmargin=4pt, aboveskip=8pt, belowskip=8pt,
  literate={é}{{\'e}}1 {è}{{\`e}}1 {ê}{{\^e}}1 {à}{{\`a}}1 {ù}{{\`u}}1
           {ô}{{\^o}}1 {î}{{\^i}}1 {ç}{{\c{c}}}1 {É}{{\'E}}1
}

\pagestyle{fancy}
\fancyhf{}
\fancyhead[L]{\small TP-1 : Setup de l'Environnement}
\fancyhead[R]{\small ESP/UCAD --- M1}
\fancyfoot[C]{\thepage}
\fancyfoot[L]{\small Dr. El Hadji Bassirou TOURÉ}
\fancyfoot[R]{\small 2025--2026}
\renewcommand{\headrulewidth}{0.4pt}
\renewcommand{\footrulewidth}{0.4pt}

\title{%
\vspace{-1cm}
\textbf{\LARGE Structures de Données et}\\
\textbf{\LARGE Algorithmes Avancés}\\[0.8cm]
\textbf{\huge TP-1 --- Setup de l'Environnement}\\[0.3cm]
\textbf{\Large Docker, Java, Gradle, JUnit, Python, Git}\\[0.5cm]
\large Master 1 --- 2025--2026}
\author{Dr. El Hadji Bassirou TOURÉ\\
Département Génie Informatique\\
Université Cheikh Anta Diop de Dakar}
\date{}

\begin{document}

\maketitle
\thispagestyle{empty}

\begin{prereqbox}[Prérequis et contexte]
\textbf{Prérequis :} Un ordinateur (Windows, macOS ou Linux) avec au moins 8~Go de RAM et 10~Go d'espace disque libre.\\[4pt]
\textbf{Objectif :} Installer et vérifier l'environnement de développement standardisé du cours. Ce TP doit être complété \textbf{avant} la séance S1.\\[4pt]
\textbf{Principe :} Vous éditez vos fichiers sur votre machine avec l'éditeur de votre choix. Vous compilez et exécutez dans le conteneur Docker via le terminal.
\end{prereqbox}

\begin{objectifbox}[Objectifs concrets]
À la fin de ce TP, vous serez capable de :
\begin{enumerate}[nosep]
    \item Lancer le conteneur Docker du cours et y accéder
    \item Compiler et exécuter un programme Java dans le conteneur
    \item Créer un projet Gradle et exécuter \texttt{gradle --no-daemon build}
    \item Écrire et exécuter un test JUnit~5
    \item Générer une figure avec Python/matplotlib
    \item Utiliser les commandes Git de base
\end{enumerate}
\end{objectifbox}

\tableofcontents
\newpage

% =====================================================================
\section{Pourquoi Docker ?}
% =====================================================================

Tous les étudiants travaillent dans le même environnement, quelle que soit leur machine. Docker crée un \textbf{conteneur} --- un environnement Linux isolé qui contient tous les outils nécessaires.

\begin{center}
\renewcommand{\arraystretch}{1.3}
\begin{tabular}{@{}l l@{}}
\toprule
\textbf{Outil} & \textbf{Usage} \\
\midrule
OpenJDK 17 (Temurin) & Compilation et exécution Java \\
Gradle 8.12 & Construction de projets, dépendances \\
JUnit 5 & Tests unitaires \\
Python 3 & Génération de figures (matplotlib) \\
Git & Contrôle de version \\
\bottomrule
\end{tabular}
\end{center}

Le principe est simple : votre dossier de travail est \textbf{partagé} entre votre machine et le conteneur. Vous éditez sur votre machine (avec l'éditeur de votre choix : VS Code, Sublime Text, Notepad++...), et vous compilez/testez dans le conteneur.

% =====================================================================
\section{Installation de Docker}
% =====================================================================

\subsection{Windows}

\begin{enumerate}[nosep]
    \item Télécharger \textbf{Docker Desktop}\footnote{\url{https://www.docker.com/products/docker-desktop/}}
    \item Lancer l'installateur (activer WSL~2 si demandé), redémarrer
    \item Installer \textbf{Git for Windows} depuis \url{https://gitforwindows.org/}
    \item Ouvrir \textbf{Git Bash} et vérifier :
\end{enumerate}

\begin{lstlisting}[style=bash]
docker --version
\end{lstlisting}

\begin{attentionbox}[Important --- Windows]
Sur Windows, utilisez toujours \textbf{Git Bash} pour les commandes Docker de ce cours (pas CMD, pas PowerShell).
\end{attentionbox}

\subsection{macOS}

\begin{enumerate}[nosep]
    \item Télécharger \textbf{Docker Desktop}\footnotemark[\value{footnote}]
    \item Choisir Apple Silicon ou Intel, installer, lancer
    \item Ouvrir Terminal et vérifier : \texttt{docker --version}
\end{enumerate}

\subsection{Linux (Ubuntu/Debian)}

\begin{lstlisting}[style=bash]
sudo apt-get update
sudo apt-get install docker.io
sudo usermod -aG docker $USER
# Se déconnecter et se reconnecter, puis vérifier :
docker --version
\end{lstlisting}

\begin{exercicebox}[Vérification 1]
La commande \texttt{docker --version} affiche un numéro de version.
\end{exercicebox}

% =====================================================================
\section{Premier lancement du conteneur}
% =====================================================================

\subsection{Créer un dossier de travail}

Créez un dossier sur votre machine pour tout le cours :

\begin{lstlisting}[style=bash]
mkdir ~/dsa-m1
cd ~/dsa-m1
\end{lstlisting}

\subsection{Lancer le conteneur}

\begin{lstlisting}[style=bash]
# Linux / macOS
docker run -it -v "$(pwd):/workspace" elbachirtoure/dsa-java-m1:v1.1

# Windows (Git Bash)
MSYS_NO_PATHCONV=1 docker run -it -v "$(pwd):/workspace" \
    elbachirtoure/dsa-java-m1:v1.1
\end{lstlisting}

La première fois, Docker télécharge l'image ($\sim$800~Mo). Les fois suivantes, le lancement est instantané.

L'option \texttt{-v} monte votre dossier local dans \texttt{/workspace}. Tout fichier créé d'un côté est visible de l'autre.

\begin{infobox}[Message d'accueil]
Vous voyez le message d'accueil du cours. Tapez \texttt{help-tp} pour l'aide-mémoire.
\end{infobox}

\subsection{Vérifier l'environnement}

\begin{lstlisting}[style=bash]
cd /workspace
java --version
gradle --version
python3 --version
git --version
\end{lstlisting}

\begin{exercicebox}[Vérification 2]
Java~17, Gradle~8.12, Python~3 et Git sont installés.
\end{exercicebox}

% =====================================================================
\section{Hello World en Java}
% =====================================================================

Sur votre machine, ouvrez votre éditeur préféré et créez le fichier \texttt{Hello.java} dans le dossier \texttt{dsa-m1/} :

\begin{lstlisting}[style=java]
public class Hello {
    public static void main(String[] args) {
        System.out.println("Structures de Donnees - M1");
        System.out.println("Environnement pret !");
    }
}
\end{lstlisting}

Dans le terminal du conteneur :

\begin{lstlisting}[style=bash]
cd /workspace
javac Hello.java
java Hello
\end{lstlisting}

\begin{exercicebox}[Vérification 3]
Le programme affiche les deux lignes. Modifiez-le pour afficher votre nom, recompilez et exécutez.
\end{exercicebox}

% =====================================================================
\section{Premier projet Gradle}
% =====================================================================

Gradle gère la compilation, les dépendances (JUnit) et les tests.

\subsection{Créer un projet}

\begin{lstlisting}[style=bash]
cd /workspace
new-project hello-dsa
cd hello-dsa
\end{lstlisting}

Structure créée :

\begin{lstlisting}[style=bash]
hello-dsa/
  src/main/java/       # Code source
  src/test/java/       # Tests JUnit
  build.gradle          # Configuration
\end{lstlisting}

\subsection{Ajouter une classe}

Sur votre machine, créez \texttt{hello-dsa/src/main/java/Main.java} :

\begin{lstlisting}[style=java]
public class Main {
    private String nom;

    public Main(String nom) {
        this.nom = nom;
    }

    public String message() {
        return "Bienvenue en DSA, " + nom + " !";
    }

    public static void main(String[] args) {
        Main s = new Main("Fatou");
        System.out.println(s.message());
    }
}
\end{lstlisting}

\subsection{Compiler et exécuter}

\begin{lstlisting}[style=bash]
gradle --no-daemon build
gradle --no-daemon run
\end{lstlisting}

\begin{infobox}[Pourquoi \texttt{-{}-no-daemon} ?]
Sur les volumes Docker partagés avec Windows, le daemon Gradle peut provoquer des erreurs d'entrée/sortie. L'option \texttt{-{}-no-daemon} évite ce problème. Le premier build prend 10--20 secondes, les suivants seront plus rapides.
\end{infobox}

\begin{exercicebox}[Vérification 4]
\texttt{gradle --no-daemon run} affiche \texttt{Bienvenue en DSA, Fatou !}
\end{exercicebox}

% =====================================================================
\section{Premier test JUnit}
% =====================================================================

Un test vérifie qu'une méthode fait ce qu'elle est censée faire.

\subsection{Écrire le test}

Créez \texttt{hello-dsa/src/test/java/MainTest.java} :

\begin{lstlisting}[style=java]
import org.junit.jupiter.api.Test;
import static org.junit.jupiter.api.Assertions.*;

class MainTest {

    @Test
    void testMessage() {
        Main s = new Main("Moussa");
        assertEquals("Bienvenue en DSA, Moussa !", s.message());
    }

    @Test
    void testAutreNom() {
        Main s = new Main("Aminata");
        assertTrue(s.message().contains("Aminata"));
    }

    @Test
    void testNonVide() {
        Main s = new Main("Cheikh");
        assertFalse(s.message().isEmpty());
    }
}
\end{lstlisting}

\subsection{Exécuter les tests}

\begin{lstlisting}[style=bash]
gradle --no-daemon test
\end{lstlisting}

Résultat attendu :

\begin{lstlisting}[style=bash]
MainTest > testMessage() PASSED
MainTest > testAutreNom() PASSED
MainTest > testNonVide() PASSED

BUILD SUCCESSFUL
\end{lstlisting}

\begin{infobox}[JUnit --- les assertions essentielles]
\begin{itemize}[nosep]
    \item \texttt{assertEquals(attendu, obtenu)} --- vérifie l'égalité
    \item \texttt{assertTrue(condition)} --- vérifie qu'une condition est vraie
    \item \texttt{assertFalse(condition)} --- vérifie qu'une condition est fausse
    \item \texttt{assertThrows(Exception.class, () -> ...)} --- vérifie qu'une exception est levée
\end{itemize}
\end{infobox}

\begin{exercicebox}[Vérification 5]
Les 3 tests passent. Ajoutez un 4e test de votre choix et relancez \texttt{gradle --no-daemon test}.
\end{exercicebox}

% =====================================================================
\section{Vérification Python et matplotlib}
% =====================================================================

Python génère les figures du cours. Créez \texttt{test\_figure.py} dans \texttt{dsa-m1/} :

\begin{lstlisting}[style=python]
import matplotlib
matplotlib.use('Agg')
import matplotlib.pyplot as plt
import numpy as np

n = np.arange(1, 50)
fig, ax = plt.subplots(figsize=(8, 5))
ax.plot(n, np.ones_like(n), label='O(1)', linewidth=2)
ax.plot(n, np.log2(n), label='O(log n)', linewidth=2)
ax.plot(n, n, label='O(n)', linewidth=2)
ax.plot(n, n * np.log2(n), label='O(n log n)', linewidth=2)

ax.set_xlabel('n')
ax.set_ylabel("Operations")
ax.set_title('Comparaison des complexites')
ax.legend()
ax.set_ylim(0, 200)
ax.grid(True, alpha=0.3)
plt.tight_layout()
fig.savefig('test_complexites.png', dpi=150)
print("Figure sauvegardee : test_complexites.png")
\end{lstlisting}

\begin{lstlisting}[style=bash]
cd /workspace
python3 test_figure.py
\end{lstlisting}

Le fichier \texttt{test\_complexites.png} apparaît dans votre dossier \texttt{dsa-m1/} --- ouvrez-le sur votre machine.

\begin{exercicebox}[Vérification 6]
La figure affiche 4 courbes de complexité correctement étiquetées.
\end{exercicebox}

% =====================================================================
\section{Git --- les bases}
% =====================================================================

\begin{lstlisting}[style=bash]
git config --global user.name "Votre Nom"
git config --global user.email "votre.email@ucad.edu.sn"

cd /workspace/hello-dsa
git init
git add .
git commit -m "TP-1 : premier projet avec JUnit"
git log --oneline
\end{lstlisting}

\begin{infobox}[Branche par défaut]
Le conteneur est configuré pour que \texttt{git init} crée une branche \texttt{main} (et non \texttt{master}).
\end{infobox}

Modifiez un fichier, puis :

\begin{lstlisting}[style=bash]
git add .
git commit -m "TP-1 : ajout d'un test"
git log --oneline
\end{lstlisting}

\begin{exercicebox}[Vérification 7]
\texttt{git log --oneline} affiche 2 commits.
\end{exercicebox}

% =====================================================================
\section{Aide-mémoire}
% =====================================================================

\begin{retenirbox}[Commandes Docker]
\renewcommand{\arraystretch}{1.3}
\small
\begin{tabular}{@{}p{8.5cm} p{5cm}@{}}
\toprule
\textbf{Commande} & \textbf{Description} \\
\midrule
\texttt{docker run -it -v "\$(pwd):/workspace" elbachirtoure/dsa-java-m1:v1.1} & Lancer le conteneur \\
\texttt{exit} & Quitter le conteneur \\
\texttt{help-tp} & Aide-mémoire (dans le conteneur) \\
\texttt{new-project <nom>} & Créer un projet Gradle \\
\bottomrule
\end{tabular}
\end{retenirbox}

\begin{retenirbox}[Outils dans le conteneur]
\renewcommand{\arraystretch}{1.3}
\begin{tabular}{@{}l l l@{}}
\toprule
\textbf{Outil} & \textbf{Commande} & \textbf{Usage} \\
\midrule
Java & \texttt{javac Fichier.java} & Compiler \\
     & \texttt{java Fichier} & Exécuter \\
Gradle & \texttt{gradle --no-daemon build} & Compiler + tests \\
       & \texttt{gradle --no-daemon test} & Tests uniquement \\
       & \texttt{gradle --no-daemon run} & Exécuter le main \\
Python & \texttt{python3 script.py} & Exécuter un script \\
Git & \texttt{git add . \&\& git commit -m "msg"} & Committer \\
\bottomrule
\end{tabular}
\end{retenirbox}

% =====================================================================
\section*{Rapport individuel}
% =====================================================================

\begin{attentionbox}[Rapport individuel --- obligatoire]
Chaque étudiant rédige un rapport de \textbf{1 à 2 pages} :
\begin{enumerate}[nosep]
    \item \textbf{Ce que j'ai fait :} outils installés, programmes écrits
    \item \textbf{Les difficultés :} erreurs rencontrées
    \item \textbf{Comment je les ai résolues}
    \item \textbf{Captures d'écran :}
    \begin{enumerate}[nosep, label=(\alph*)]
        \item Message d'accueil du conteneur
        \item Résultat de \texttt{gradle --no-daemon test} avec les tests passés
        \item La figure \texttt{test\_complexites.png}
    \end{enumerate}
\end{enumerate}
\textbf{Format :} PDF. \textbf{Nom :} \texttt{rapport-tp1-PRENOM-NOM.pdf}\\
\textbf{Remise :} Avant la séance S1.
\end{attentionbox}

\vfill
\begin{center}
\rule{0.6\textwidth}{0.4pt}\\[6pt]
{\small Dr. El Hadji Bassirou TOURÉ --- ESP/UCAD --- M1 --- 2025--2026}
\end{center}

\end{document}
